% !BIB program = biber
% 中文简历入口；与 main.tex 共用 res.cls 和 mypub.bib。
% 同时支持 XeLaTeX 和 pdfLaTeX，文件保持 UTF-8 编码。
% 编译：所选引擎 main_chinese.tex -> biber main_chinese
%       -> 所选引擎 main_chinese.tex -> 所选引擎 main_chinese.tex
% 自动编译：latexmk -xelatex main_chinese.tex 或 latexmk -pdf main_chinese.tex
% Overleaf：主文档选择 main_chinese.tex，编译器可选 XeLaTeX 或 pdfLaTeX。
% 切换编译器时，请清理辅助文件（Overleaf：Recompile from scratch）。
% pdfLaTeX 需要 CJK/ctex 和 Arphic 字体包（含 gbsn 宋体）。
\documentclass[margin]{res}

% scheme=plain 保留 res.cls 的章节样式；按引擎选择中文字体。
% 两种配置均使用 TeX 发行版提供的字体，不依赖系统专有中文字体。
\usepackage{iftex}
\ifPDFTeX
  \usepackage[T1]{fontenc}
  \usepackage{lmodern}
  % ctex 自动启用 CJKutf8；禁用不兼容 pdfLaTeX 的 Fandol 字体配置。
  \usepackage[UTF8,scheme=plain,fontset=none]{ctex}
  \renewcommand{\CJKrmdefault}{gbsn}
  \renewcommand{\CJKsfdefault}{gbsn}
  \renewcommand{\CJKttdefault}{gbsn}
\else
  \usepackage[UTF8,scheme=plain,fontset=fandol]{ctex}
\fi
\usepackage{helvet}
\usepackage{xcolor}

% 用现代字体命令替代旧模板的 \bf，保持中英文标题的粗体一致。
\renewcommand{\sectionfont}{\bfseries}
\renewcommand{\namefont}{\large\bfseries}

\textheight=700pt

% res.cls 默认使用 \nofiles；恢复辅助文件写入以支持 biblatex/Biber。
\makeatletter
\@fileswtrue
\makeatother
\usepackage[backend=biber,style=numeric,sorting=ydnt,defernumbers=true,
  giveninits=true,maxbibnames=99,doi=true,url=true,isbn=false]{biblatex}
\addbibresource{mypub.bib}

% 与英文版一致：仅强调 Yuanchao Xu，兼容全名和姓名缩写。
\renewcommand*{\mkbibcompletename}[1]{%
  \ifboolexpr{test {\ifdefstring{\namepartfamily}{Xu}} and
              (test {\ifdefstring{\namepartgiven}{Yuanchao}} or test {\ifdefstring{\namepartgiven}{Y.}} or test {\ifdefstring{\namepartgiven}{Y}})}
    {\textcolor{black}{\underline{\textbf{Y.\nobreakspace Xu}}}}
    {#1}}
\setlength{\bibitemsep}{0.5\baselineskip}
\setlength{\emergencystretch}{2em}

% 仅在中文版中翻译文献状态；论文题名、作者及出版信息保持原文。
\DeclareFieldFormat{note}{%
  \ifstrequal{#1}{Preprint}{预印本}{%
  \ifstrequal{#1}{Poster}{海报}{%
  \ifstrequal{#1}{Submitted}{已投稿}{%
  \ifstrequal{#1}{In preparation}{准备中}{#1}}}}}

\usepackage[hidelinks,unicode]{hyperref}
\hypersetup{pdftitle={Yuanchao Xu - 中文简历},pdfauthor={Yuanchao Xu}}
\renewcommand{\today}{\number\year 年\number\month 月\number\day 日}

\begin{document}

% 原项目只提供英文姓名；可在此替换为本人确认的中文姓名。
\name{\href{https://sites.google.com/view/yuanchao/home}{Yuanchao Xu}}
\address{%
  \begin{minipage}{\linewidth}
    \centering
    \href{mailto:xu.yuanchao.3a@kyoto-u.ac.jp}{xu.yuanchao.3a@kyoto-u.ac.jp}\\[0.5ex]
    京都大学 益川教育研究馆 4-405 室\\
    \today
  \end{minipage}%
}

\begin{resume}

\section{工作经历}

\textbf{阿尔伯塔大学}\\
访问博士后，应用数学\\
2026年4月 -- 2026年8月

\textbf{京都大学}\\
特定研究员，\href{https://sci.kyoto-u.ac.jp/en/search/tag?tag=SACRA}{SACRA}\\
2025年9月 -- 2030年8月

\section{教育经历}

\textbf{阿尔伯塔大学}\\
应用数学博士（获博士学位论文奖）\\
2021年9月 -- 2025年8月

\textbf{曼尼托巴大学}\\
数学理学硕士\\
2018年9月 -- 2020年8月

\textbf{俄勒冈大学}\\
数学理学学士（荣誉毕业，Cum Laude）\\
2015年1月 -- 2017年12月

\section{研究兴趣}

我的研究聚焦于随机动力系统的数据驱动方法，尤其关注学习随机 Koopman 算子的理论框架，以及这些方法在计算神经科学、强化学习和生成建模等复杂实际问题中的应用。

\begin{format}
\title{l}\\
\dates{l}\location{r}\\
\body\\
\end{format}

\section{论文与预印本}
\nocite{*}
\printbibliography[heading=none,notkeyword=duplicate,notkeyword=draft]

\section{已投稿与\\准备中论文}
\printbibliography[heading=none,keyword=draft]

\section{学术报告}
\begin{enumerate}
    \item 基于 Koopman 谱分析的生成建模\\
    \textbf{麻省理工学院}，美国马萨诸塞州剑桥，2026年4月27日

    \item 基于 Koopman 谱分析的生成建模\\
    CMS 专题研讨会\\
    \textbf{加州理工学院}，美国加利福尼亚州帕萨迪纳，2026年4月13日

    \item 基于 Koopman 谱分析的生成建模\\
    \href{https://naturalsciences.uoregon.edu/mathematics/seminars}{概率论研讨会}\\
    \textbf{俄勒冈大学}，美国俄勒冈州尤金，2026年4月8日

    \item 基于 Koopman 谱分析的生成建模\\
    \href{https://sites.google.com/view/mlds-deds-2026}{MLDS 5 \& DEDS 2026}\\
    \textbf{京都大学}，日本京都，2026年2月9--13日

    \item 基于 Koopman 谱分析的生成建模\\
    \href{https://www.youtube.com/watch?v=FDEGkddluls&t=893s}{机器学习与动力系统研讨会}\\
    \textbf{艾伦·图灵研究所}，英国伦敦，2026年1月8日（线上）

    \item 复杂动力系统建模的数据驱动 Koopman 框架\\
    \textbf{复旦大学}，中国上海，2026年1月5日

    \item 复杂动力系统建模的数据驱动 Koopman 框架\\
    \textbf{欧洲分子生物学实验室}，意大利罗马，2025年11月24日（线上）

    \item 复杂动力系统建模的数据驱动 Koopman 框架\\
    \textbf{于贝尔·居里安实验室（Laboratoire Hubert Curien）}，法国圣艾蒂安，2025年10月14日

    \item 学习随机 Koopman 算子的扰动方法\\
    \href{https://docs.google.com/spreadsheets/d/1XrgG53J6HlsHYLRHtYFsj5s32eWDr9B_/edit?gid=1512777206#gid=1512777206}{CREST}\\
    \textbf{理化学研究所（RIKEN）}，日本神户，2024年12月26--28日

    \item 从复杂系统中提取动力学：确定性与随机性视角\\
    \textbf{华东理工大学}，中国上海，2024年12月19日
\end{enumerate}

\section{学术访问}
% 原英文文件前五项未标年份，此处保留原始日期，不推测补全。
\begin{enumerate}
    \item 麻省理工学院，美国马萨诸塞州剑桥，4月27--30日
    \item 康奈尔大学，美国纽约州伊萨卡，4月16--18日
    \item 加州理工学院，美国加利福尼亚州帕萨迪纳，4月13--15日
    \item 加州大学圣巴巴拉分校，美国加利福尼亚州圣巴巴拉，4月10日
    \item 俄勒冈大学，美国俄勒冈州尤金，4月6--9日
    \item 复旦大学，中国上海，2026年1月5--9日
    \item 于贝尔·居里安实验室（Laboratoire Hubert Curien）\\
    法国圣艾蒂安，2025年10月14--15日
    \item 欧洲分子生物学实验室，意大利罗马，2025年10月11--13日
    \item 意大利技术研究院（Istituto Italiano di Tecnologia）\\
    意大利热那亚，2025年10月9--10日
    \item 爱媛大学，日本松山，2024年12月23--25日
\end{enumerate}

\section{审稿工作}

ICML 2026、ICLR 2026、IEEE-TNNLS、IEEE-TETCI

\section{学术服务}

DCDS-I 青年编委（Young Associate Editor）

\section{奖励与荣誉}

\textbullet{} 理学院博士学位论文奖（Faculty of Science Doctoral Dissertation Award），2025年\\
\textbullet{} Mary Louise Imrie 研究生奖，2025年\\
\textbullet{} Dr. Josephine M. Mitchell 招生奖学金，2021年\\
\textbullet{} 国际研究生入学奖学金，2020年\\
\textbullet{} 国际研究生入学奖学金，2018年\\
\textbullet{} Clarence and Lucille Dunbar 奖学金，2017年

\section{教学经历}

Math 102 应用线性代数（习题课）\hfill 2025年冬季\\
Math 209 工程微积分 III（习题课）\hfill 2024年秋季\\
Math 209 工程微积分 III（习题课）\hfill 2024年春季\\
Math 102 应用线性代数（习题课）\hfill 2024年冬季\\
Math 201 微分方程（习题课）\hfill 2023年秋季\\
Math 101 工程微积分 II（习题课）\hfill 2023年春季\\
Math 201 微分方程（习题课）\hfill 2023年冬季\\
研究生科研助理\hfill 2020年9月 -- 2021年8月\\
研究生教学助理\hfill 2018年9月 -- 2020年8月\\
本科生教学助理\hfill 2016年9月 -- 2017年12月\\
本科生数学辅导教师\hfill 2016年4月 -- 2017年12月

\end{resume}
\end{document}
